\SeventhArticle{16}
Die ganze Aufgabe verlangt von uns, für jeden gegebenen Wert von $n$ eine Funktion $T$ der Form $t^{n+1}+\alpha t^n+\alpha't^{n-1}+\alpha''t^{n-2}$ usw. zu finden, die so beschaffen ist, dass bei der Entwicklung des Produkts
\[
T(t^{-1}+\tfrac12t^{-2}+\tfrac13t^{-3}+\tfrac14t^{-4}+\text{usw.})
\]
die Potenzen $t^{-1},t^{-2},t^{-3}\mathinner{.\,.\,.\,.}t^{-(n+1)}$ sämtlich den Koeffizienten $0$ erhalten; oder, wenn man es vorzieht, eine Funktion $U$ der Form $u^{n+1}+\beta u^n+\beta'u^{n-1}+\beta''u^{n-2}+$ usw., deren Produkt mit $u^{-1}+\frac13u^{-3}+\frac15u^{-5}+\frac17u^{-7}+$ usw. die Potenzen \SourceLeaf{193}%
$u^{-1},u^{-2},u^{-3},u^{-4}\mathinner{.\,.\,.\,.}u^{-(n+1)}$ nicht enthält. Das letztere Verfahren wird etwas einfacher sein: Denn man sieht leicht, dass die Koeffizienten von $U$ abwechselnd verschwinden müssen, damit die vorgeschriebene Bedingung erfüllt ist, oder dass man $\beta=0,\beta''=0,\beta''''=0$ usw. setzen muss; damit kann beinahe die Hälfte der Arbeit bereits als erledigt gelten. Entwickeln wir einige einfachere Fälle.

I. Für $n=0$ muss der einzige Koeffizient von $t^{-1}$ im Produkt
\[
(t+\alpha)(t^{-1}+\tfrac12t^{-2}+\tfrac13t^{-3}+\text{usw.})
\]
verschwinden. Da er $=\frac12+\alpha$ wird, haben wir $\alpha=-\frac12$ oder $T=t-\frac12$. Ebenso ist $U=u$.

II. Für $n=1$ hängt die Bestimmung von $T$ von zwei Gleichungen ab:
\[
\begin{aligned}
0&=\tfrac13+\tfrac12\alpha+\alpha'\\
0&=\tfrac14+\tfrac13\alpha+\tfrac12\alpha'
\end{aligned}
\]
Daraus leiten wir $\alpha=-1,\alpha'=+\frac16$ oder $T=tt-t+\frac16$ ab. Die Bestimmung der Funktion $U$ ergibt nur eine Gleichung:
\[
0=\tfrac13+\beta'
\]
Daraus folgt $\beta'=-\frac13$ oder $U=uu-\frac13$.

III. Für $n=2$ wird die Funktion $T$ mithilfe dreier Gleichungen bestimmt:
\[
\begin{aligned}
0&=\tfrac14+\tfrac13\alpha+\tfrac12\alpha'+\alpha''\\
0&=\tfrac15+\tfrac14\alpha+\tfrac13\alpha'+\tfrac12\alpha''\\
0&=\tfrac16+\tfrac15\alpha+\tfrac14\alpha'+\tfrac13\alpha''
\end{aligned}
\]
Daraus erhalten wir $\alpha=-\frac32,\alpha'=\frac35,\alpha''=-\frac1{20}$ und somit $T=t^3-\frac32tt+\frac35t-\frac1{20}$. Zur Bestimmung von $U$ genügt eine einzige Gleichung:
\[
0=\tfrac15+\tfrac13\beta'
\]
Daraus folgt $\beta'=-\frac35$ oder $U=u^3-\frac35u$.

Wir werden dieses Verfahren, das zu immer mühsameren Rechnungen führt, hier jedoch nicht weiter verfolgen, sondern zum eigentlichen Ursprung der allgemeinen Lösung übergehen.
\SourceLeaf{194}
\SeventhArticle{17}
Für einen gegebenen Kettenbruch
\[
\varphi=\cfrac{v}{w+\cfrac{v'}{w'+\cfrac{v''}{w''+\cfrac{v'''}{w'''+\text{usw.}}}}}
\]
findet man bekanntlich immer näher kommende Näherungsbrüche durch den folgenden Algorithmus. Man bilde zwei Reihen von Größen, $V,V',V'',V'''$ usw., $W,W',W'',W'''$ usw., nach diesen Formeln:
\[
\begin{array}{ll}
V=0&W=1\\[4pt]
V'=v&W'=wW\\[4pt]
V''=w'V'+v'V&W''=w'W'+v'W\\[4pt]
V'''=w''V''+v''V'&W'''=w''W''+v''W'\\[4pt]
V''''=w'''V'''+v'''V''&W''''=w'''W'''+v'''W''
\end{array}
\]
usw.; dann ist
\[
\begin{aligned}
\cfrac VW&=0\\[9pt]
\cfrac{V'}{W'}&=\cfrac vw\\[9pt]
\cfrac{V''}{W''}&=\cfrac v{w+\cfrac{v'}{w'}}\\[9pt]
\cfrac{V'''}{W'''}&=\cfrac v{w+\cfrac{v'}{w'+\cfrac{v''}{w''}}}
\end{aligned}
\]
und so weiter. Ferner ist bekannt oder lässt sich aus den vorstehenden Gleichungen selbst leicht bestätigen, dass gilt:
\[
\begin{aligned}
VW'-V'W&=-v\\
V'W''-V''W'&=+vv'\\
V''W'''-V'''W''&=-vv'v''\\
V'''W''''-V''''W'''&=+vv'v''v'''
\end{aligned}
\]
usw. Daher gilt für die Reihe
\[
\frac v{WW'}-\frac{vv'}{W'W''}+\frac{vv'v''}{W''W'''}-\frac{vv'v''v'''}{W'''W''''}+\text{usw.}
\]
\SourceLeaf{195}
\[
\begin{array}{ll}
\text{erstes Glied}&=\dfrac{V'}{W'}\\[10pt]
\text{Summe der ersten beiden Glieder}&=\dfrac{V''}{W''}\\[10pt]
\text{Summe der ersten drei Glieder}&=\dfrac{V'''}{W'''}\\[10pt]
\text{Summe der ersten vier Glieder}&=\dfrac{V''''}{W''''}
\end{array}
\]
und so weiter; folglich wird die Reihe selbst, bis ins Unendliche oder bis zu ihrem Abbruch fortgesetzt, den Kettenbruch $\varphi$ selbst ausdrücken. Zugleich ergibt sich daraus die Differenz zwischen $\varphi$ und den einzelnen Näherungsbrüchen $\frac{V'}{W'},\frac{V''}{W''},\frac{V'''}{W'''}$ usw.

Aus Formel 33 des Artikels 14 der \textit{Allgemeinen Untersuchungen über die unendliche Reihe} erhalten wir, indem wir $t$ durch $\frac1u$ ersetzen, leicht die Umwandlung der Reihe
\[
\varphi=u^{-1}+\tfrac13u^{-3}+\tfrac15u^{-5}+\tfrac17u^{-7}+\text{usw.}
\]
in den folgenden Kettenbruch,
\[
\cfrac{1}{u-\cfrac{\cfrac13}{u-\cfrac{\cfrac{2\mathbin{.}2}{3\mathbin{.}5}}{u-\cfrac{\cfrac{3\mathbin{.}3}{5\mathbin{.}7}}{u-\cfrac{\cfrac{4\mathbin{.}4}{7\mathbin{.}9}}{u-\text{usw.}}}}}}
\]
sodass gilt:
\[
\begin{aligned}
v&=1,\ v'=-\tfrac13,\ v''=-\tfrac4{15},\ v'''=-\tfrac9{35},\ v''''=-\tfrac{16}{63}\ \text{usw.}\\
w&=w'=w''=w'''=w''''\ \text{usw.}=u.
\end{aligned}
\]
Daraus erhalten wir für $V,V',V'',V'''$ usw., $W,W',W'',W'''$ usw. die folgenden Werte:
\[
\begin{aligned}&V=0,\\[5pt]&\qquad W=1\\[5pt]&V'=1,\\[5pt]&\qquad W'=u\\[5pt]&V''=u,\\[5pt]&\qquad W''=uu-\tfrac13\\[5pt]&V'''=uu-\tfrac4{15},\\[5pt]&\qquad W'''=u^3-\tfrac35u\end{aligned}\]
\[\begin{aligned}&V''''=u^3-\tfrac{11}{21}u,\\[5pt]&\qquad W''''=u^4-\tfrac67uu+\tfrac3{35}\\[5pt]&V^{\mathrm{V}}=u^4-\tfrac79uu+\tfrac{64}{945},\\[5pt]&\qquad W^{\mathrm{V}}=u^5-\tfrac{10}{9}u^3+\tfrac5{21}u\\[5pt]&V^{\mathrm{VI}}=u^5-\tfrac{34}{33}u^3+\tfrac15u,\\[5pt]&\qquad W^{\mathrm{VI}}=u^6-\tfrac{15}{11}u^4+\tfrac5{11}uu-\tfrac5{231}\\[5pt]&V^{\mathrm{VII}}=u^6-\tfrac{50}{39}u^4+\tfrac{283}{715}uu-\tfrac{256}{15015},\\[5pt]&\qquad W^{\mathrm{VII}}=u^7-\tfrac{21}{13}u^5+\tfrac{105}{143}u^3-\tfrac{35}{429}u\ \text{usw.}\end{aligned}
\]
\SourceLeaf{196}
Mit etwas Aufmerksamkeit wird klar, dass die einzelnen Größen $V,V',V'',V'''$ usw., $W,W',W'',W'''$ usw. Polynome in der Unbestimmten $u$ sind; das höchste Glied in $V^{(m)}$ ist $u^{m-1}$, und die Potenzen $u^{m-2},u^{m-4},u^{m-6}$ usw. fehlen; das höchste Glied in $W^{(m)}$ ist hingegen $u^m$, und die Potenzen $u^{m-1},u^{m-3},u^{m-5}$ usw. fehlen. Nach dem, was wir oben bewiesen haben, ist aber
\[
\begin{aligned}&\varphi=\frac1{WW'}+\frac1{3W'W''}+\frac{2\mathbin{.}2}{3\mathbin{.}3\mathbin{.}5W''W'''}\\[5pt]&+\frac{2\mathbin{.}2\mathbin{.}3\mathbin{.}3}{3\mathbin{.}3\mathbin{.}5\mathbin{.}5\mathbin{.}7W'''W''''}+\frac{2\mathbin{.}2\mathbin{.}3\mathbin{.}3\mathbin{.}4\mathbin{.}4}{3\mathbin{.}3\mathbin{.}5\mathbin{.}5\mathbin{.}7\mathbin{.}7\mathbin{.}9W''''W'''''}+\text{usw.}\end{aligned}
\]
und somit allgemein
\[
\begin{aligned}
\varphi-\frac{V^{(m)}}{W^{(m)}}&=\frac{2\mathbin{.}2\mathbin{.}3\mathbin{.}3\mathinner{.\,.\,.\,.\,.} m\mathbin{.}m}{3\mathbin{.}3\mathbin{.}5\mathbin{.}5\mathinner{.\,.\,.\,.}(2m-1)(2m+1)W^{(m)}W^{(m+1)}}\\[5pt]
&\quad+\frac{2\mathbin{.}2\mathbin{.}3\mathbin{.}3\mathinner{.\,.\,.\,.\,.}(m+1)(m+1)}{3\mathbin{.}3\mathbin{.}5\mathbin{.}5\mathinner{.\,.\,.\,.\,.}(2m+1)(2m+3)W^{(m+1)}W^{(m+2)}}\\[5pt]
&\quad+\text{usw.}
\end{aligned}
\]
Wird daher $\varphi-\frac{V^{(m)}}{W^{(m)}}$ in eine absteigende Reihe entwickelt, so lautet ihr erstes Glied
\[
=\frac{2\mathbin{.}2\mathbin{.}3\mathbin{.}3\mathinner{.\,.\,.\,.} m\mathbin{.}m u^{-(2m+1)}}{3\mathbin{.}3\mathbin{.}5\mathbin{.}5\mathinner{.\,.\,.\,.}(2m-1)(2m+1)}
\]
Das Produkt $\varphi W^{(m)}$ besteht hingegen aus dem Polynom $V^{(m)}$ und einer unendlichen Reihe, deren erstes Glied lautet
\[
=\frac{2\mathbin{.}2\mathbin{.}3\mathbin{.}3\mathinner{.\,.\,.\,.\,.} mmu^{-(m+1)}}{3\mathbin{.}3\mathbin{.}5\mathbin{.}5\mathinner{.\,.\,.\,.\,.}(2m-1)(2m+1)}
\]
Damit ist die Funktion $U$ vom Grad $n+1$, welche die im vorigen Artikel festgelegte Bedingung erfüllt, von selbst gefunden: Das Produkt $\varphi U$ soll nämlich die Potenzen $u^{-1},u^{-2},u^{-3}\mathinner{.\,.\,.\,.\,.} u^{-(n+1)}$ nicht enthalten. Sie ist keine andere als $W^{(n+1)}$; zugleich ist klar, dass $U'$ gleich $V^{(m+1)}$ ist und das erste Glied von $U''$ lautet
\[
=\frac{2\mathbin{.}2\mathbin{.}3\mathbin{.}3\mathinner{.\,.\,.\,.\,.}(n+1)(n+1)}{3\mathbin{.}3\mathbin{.}5\mathbin{.}5\mathinner{.\,.\,.\,.\,.}(2n+1)(2n+3)}\mathbin{.}u^{-(n+2)}
\]
Nimmt man daher für $b,b',b''\mathinner{.\,.\,.\,.\,.} b^{(n)}$ die Wurzeln der Gleichung $W^{(n+1)}=0$ und ermittelt die Werte der Koeffizienten $R,R',R''\mathinner{.\,.\,.\,.\,.} R^{(n)}$ nach den oben gegebenen Regeln, so besitzt unsere Integralformel eine Genauigkeit, die bis zur Ordnung $2n+1$ reicht, und ihre Korrektur wird näherungsweise ausgedrückt durch
\[
\begin{aligned}&\frac1{2^{2n+2}}\mathbin{.}\frac{2\mathbin{.}2\mathbin{.}3\mathbin{.}3\mathinner{.\,.\,.\,.\,.}(n+1)(n+1)}{3\mathbin{.}3\mathbin{.}5\mathbin{.}5\mathinner{.\,.\,.\,.\,.}(2n+1)(2n+3)}L^{(2n+2)}\\[5pt]&=\frac{1\mathbin{.}1\mathbin{.}2\mathbin{.}2\mathbin{.}3\mathbin{.}3\mathinner{.\,.\,.\,.\,.}(n+1)(n+1)}{2\mathbin{.}6\mathbin{.}6\mathbin{.}10\mathbin{.}10\mathbin{.}14\mathinner{.\,.\,.\,.\,.}(4n+2)(4n+6)}L^{(2n+2)}\end{aligned}
\]
\SourceLeaf{197}
